\documentclass{beamer}

\setbeamertemplate{bibliography item}{[\theenumiv]}

\usetheme{Madrid}
%\usecolortheme{seahorse}
\usepackage[brazil]{babel}
\usepackage[utf8x]{inputenc} % acentos diretamente do teclado
\usepackage[T1]{fontenc}
\usepackage{graphicx}


\title[Programação Paralela]{Intel Cilk Plus}
\subtitle{Programação Paralela}
\author[Alexandre \and Jeremias \and Matheus]{Alexandre Lucchesi%
    \and Jeremias Moreira%
    \and Matheus Braga}
\institute[UnB]{%
    Departamento de Ciência da Computação\\
    Universidade de Brasília, Brasília -- DF\\[1ex]
    \texttt{alexandrelucchesi@gmail.com}\\
    \texttt{jeremias@aluno.unb.br}\\
    \texttt{matheus.mtb7@gmail.com}\\
}
\date[Outubro, 2014]{10 de outubro de 2014}

\begin{document}

\begin{frame}[plain]
    \titlepage%
\end{frame}

\begin{frame}[shrink]{Sumário}
    \tableofcontents
\end{frame}


\section{Introdução}
\subsection{O que é}
\begin{frame}{Introdução}
    \begin{itemize}
        \item Cilk Plus é ``dahora''~\cite{jeffers:2013}.
    \end{itemize}
\end{frame}

\subsection{História}
\begin{frame}
    É da hora mesmo\ldots
\end{frame}

\section{Principais Funções}
\subsection{Palavras-chave}
{%
\usebackgroundtemplate{\includegraphics[width=\paperwidth,height=\paperheight]{./img/spawn.jpg}}
\begin{frame}{\texttt{cilk\_spawn}}
\end{frame}
}

\subsection{Exemplos}
\begin{frame}
    \begin{block}{Teste}
        Mais teste
    \end{block}
\end{frame}


\section{Análise Comparativa}
\subsection{Gráfico de \protect\textit{Speedup}}


\section{Aspectos da Ferramenta}
\subsection{Vantagens}
\subsection{Desvantagens}


\section{Conclusão}
\begin{frame}%[allowframebreaks]
    \frametitle{Referências Bibliográficas}
    \tiny{\bibliographystyle{abbrv}}
    \bibliography{refs}
\end{frame}

\end{document}

